\newpage
\phantomsection
\label{prf:order-comparison-of-suprema}

\begin{remark*}[Return]
\hyperref[prop:order-comparison-of-suprema]{Return to Theorem}
\end{remark*}

\begin{theorem*}[Order Comparison of Supremum]
Let $A, B \subseteq \mathbb{R}$ be non-empty sets that are bounded above. If for every $x \in A$, there exists $y \in B$ such that $x \le y$, then $\sup A \le \sup B$.
\end{theorem*}

\begin{proof}
\textbf{Professional Standard Proof.}~\\
Let $s_B = \sup B$. By definition, $y \le s_B$ for all $y \in B$. By hypothesis, for any $x \in A$, there exists $y \in B$ such that $x \le y$. By transitivity, $x \le s_B$ for all $x \in A$, making $s_B$ an upper bound for $A$. Since $\sup A$ is the least upper bound of $A$, it follows that $\sup A \le s_B = \sup B$.
\end{proof}

\begin{proof}
\textbf{Detailed Learning Proof.}~\\

\textbf{Step 1. Identify the upper bound of the target set.}
Let $s_B = \sup B$. By the definition of the supremum as an upper bound, we have:
\[ \forall y \in B, y \le s_B \]

\textbf{Step 2. Relate elements of $A$ to the upper bound of $B$.}
Let $x \in A$ be arbitrary. The hypothesis states that there exists some $y \in B$ such that $x \le y$. Combining this with Step 1, we have:
\[ x \le y \quad \text{and} \quad y \le s_B \]
By the transitivity of the order relation on $\mathbb{R}$:
\[ x \le s_B \]

\textbf{Step 3. Establish that $s_B$ bounds set $A$.}
Since $x$ was chosen arbitrarily from $A$, the inequality $x \le s_B$ holds for all $x \in A$. Therefore, $s_B$ is an upper bound for the set $A$.

\textbf{Step 4. Apply the minimality of the supremum.}
By the definition of $\sup A$ as the \textit{least} upper bound, it must be less than or equal to any other upper bound of $A$. Since $s_B$ is an upper bound for $A$, we conclude:
\[ \sup A \le s_B \]
Substituting back $s_B = \sup B$, we reach the desired conclusion:
\[ \sup A \le \sup B \]
\end{proof}

\begin{remark*}[Proof structure]
The proof employs a direct argument by showing that the supremum of $B$ serves as an upper bound for $A$, then invoking the "leastness" property of the supremum of $A$.
\end{remark*}

\begin{remark*}[Dependencies]~\\
\begin{itemize}
  \item \hyperref[def:supremum]{Definition of Supremum}
  \item \hyperref[def:upper-bound]{Definition of Upper Bound}
  \item \hyperref[ax:transitivity-of-order]{Transitivity of $\le$ on $\mathbb{R}$}
\end{itemize}
\end{remark*}